\chapter{Secure Architecture Document}
\chaptermark{Secure Architecture}
\label{ch:secure-arch}

This chapter details the defense-in-depth strategy for the secure online banking architecture. It explains how the three-tier Azure infrastructure is partitioned, how data is protected in transit, and how cryptographic secrets are managed throughout their lifecycle.

\section{Network topology and firewall placement}

The banking infrastructure uses a segmented network design to isolate sensitive data from public-facing components. Traffic is strictly filtered at each transition point so that a compromise in one zone does not grant automatic access to another.

\subsection*{Public zone (Internet edge)}
The entry point for customers, tellers, and admins. Traffic is first filtered by an \emph{external firewall} that provides DDoS mitigation and rate limiting to reduce the risk of service exhaustion.

\subsection*{Demilitarized zone (DMZ)}
Houses the Web Application Firewall (WAF) and the frontend Python Flask server. The WAF inspects traffic for application-layer attacks.

\subsection*{Private application zone}
Separated by an \emph{internal firewall}, this zone contains the C\# .NET Core Web API. This layer enforces least privilege: only validated requests from the frontend reach the business logic.

\subsection*{Restricted data zone}
The most secure layer, containing the Azure SQL Database. Access is restricted by an Azure SQL firewall that allows only the specific IP addresses of the web API.

\begin{figure}[htbp]
  \centering
  \includegraphics[width=\textwidth,height=0.88\textheight,keepaspectratio]{Blank_board.pdf}
  \caption{Network topology and firewall placement (reference diagram).}
  \label{fig:network-topology-board}
\end{figure}

Figure~\ref{fig:network-topology-board} summarizes the segmented layout and trust boundaries for the banking infrastructure.

\section{TLS implementation plan}

To protect the confidentiality and integrity of banking data in transit, Transport Layer Security (TLS)~1.3 is enforced across communication channels.

\subsection*{Implementation scope}

\begin{description}
  \item[User to frontend:]
    All web traffic is encrypted via HTTPS/TLS~1.3 to protect login credentials and session tokens.
  \item[Service to service:]
    Communication between the frontend UI and backend API uses TLS~1.3 to mitigate information disclosure.
  \item[Database connectivity:]
    The system uses SQL over TLS for all database queries to reduce the risk of sniffing sensitive financial records.
\end{description}

\subsection*{Certificate and protocol hardening}

The project treats downgrade attacks on older browsers as a risk. To mitigate this, HTTP Strict Transport Security (HSTS) headers are implemented so browsers connect only via HTTPS.

\section{Key management strategy}

Encryption depends on protecting cryptographic keys. This project centralizes cryptographic operations and secret storage in Azure Key Vault.

\subsection*{Secure storage and separation}

Cryptographic keys and database credentials are stored in Key Vault, separate from application code and the database tier.

\subsection*{Access control and identity}

\begin{itemize}
  \item The system uses managed identities and role-based access control (RBAC) so that only authorized services (for example the backend API) can retrieve specific secrets.
  \item Least-privilege access policies are enforced to reduce the risk of broad credential theft.
\end{itemize}

\subsection*{Secret protection}

Secure password hashing uses bcrypt/Argon2 with cryptographically secure random salts for every user.

\subsection*{Audit integrity}

Access to keys is logged. To prevent a compromised administrator from destroying evidence, audit logs are stored on a separate server with hash-protected integrity checks to detect any alteration.
